% !TeX root = ../main.tex
% !TeX encoding = UTF-8
% !TeX spellcheck = none

\section{讨论}\label{sec:discussion}

CMSRN围绕常规流量模式偏离时的预测修正，将异质证据组织与历史消息选择纳入共享双频表征之上的残差框架。SERR按自身演化和有向时滞两类角色组织证据，BCMS依据流量背景与预测时距选择跨区域历史消息。两者从共享表征分别生成修正并作用于预测端，使不同来源的信息能够以明确的功能分工参与多步预测。

总体结果与情景分析共同说明，这一设计的预测优势具有状态和任务差异。在主要实验的13种基线范围内，CMSRN在两城市的流入、流出任务上均取得最低AVG12 MAE，相对最强基线降低0.46\%--2.00\%；其中成都流入的AVG12 RMSE降低5.01\%，上海流出的AVG12 RMSE则略高于最强基线。加入STNorm和STWave的扩展比较进一步显示，CMSRN在四项任务中的三项Extreme组取得最低MAE，而多个较低偏离组仍由上述强基线取得更低误差。这表明，整体误差、局部偏离程度与流量通道需要结合评价。空间案例中保留主要热点位置但存在幅值偏差的现象，也提示区域调度应用应同时关注流量分布与局部峰值误差。

两条残差通路的分组响应为理解这种差异提供了依据。自身修正在近期快速变化的高偏离样本中对应更大的绝对MAE收益，时滞修正提供幅度较小、随城市而变化的补充；BCMS的平均收益则在目标与候选背景代理较接近时更大，并覆盖高偏离子集之外的交通状态。这些观察支持将局部演化证据与跨区域历史信息分工处理，也表明跨区域补充的作用范围不限于强偏离情景。第\ref{sec:exp-ablation}节的独立重训练消融评价各修正通路的整体贡献，固定输出诊断则解释完整模型对既有残差的依赖，两类证据从不同角度刻画了角色修正的作用。

本研究在成都和上海各一个月的规则网格数据上，考察了CMSRN对5–60分钟交通流量的预测能力。现有实验主要反映观测时段内的短时预测表现，其在跨季节变化和持续性交通模式迁移下的适应能力值得进一步研究。当前常规参照、有向时滞关系与候选支持在预测阶段保持固定，如何随交通状态演变更新这些先验信息，是拓展模型时间适应性的关键。

未来研究将进一步探索时间适应性。面向更长时间跨度的交通变化，研究常规参照、有向时滞关系与候选支持的动态更新，使预测修正能够持续适应周期规律和跨区域关联的演变。在此基础上，结合预测目标严格隔离的评估，以及对零流量区域和局部峰值的分析，进一步检验模型在不同流量状态下的适用性，为持续变化的交通环境提供更可靠的短时预测支持。
